\clearpage\thispagestyle{empty}
\pdfbookmark[0]{Quellen zu den Elementen der Psyche}{psyche818-sources}
\begin{center}\textbf{Quellen zu den Elementen der Psyche}\end{center}
\noindent Diese KI-Übertragung ins heutige Deutsch folgt Carl Friedrich Gauß’ Bericht auf S.~479 der \href{https://gdz.sub.uni-goettingen.de/id/PPN236007467}{\textit{Werke}, Band~VI} (Göttingen, 1874). Gauß erörtert die Bahn der Psyche und gibt die von Wilhelm Klinkerfues bestimmten Elemente an. Die Abkürzung Th.\ M.\ C.\ C.\ bezeichnet Gauß’ \textit{Theoria motus corporum coelestium in sectionibus conicis solem ambientium} (Hamburg, 1809).
\par\smallskip\noindent Der Bericht erschien in den \textit{Astronomischen Nachrichten}, Band~XXXV, Nr.~818, \href{https://api.digitale-sammlungen.de/iiif/image/v2/bsb10049460_00013/full/1800,/0/default.jpg}{Sp.~17}. Die Zeitschrift enthält eine Vergleichstabelle. Gauß’ Unterschrift und das Datum 8. Juli 1852 stehen am Ende des Beitrags auf \href{https://api.digitale-sammlungen.de/iiif/image/v2/bsb10049460_00014/full/1800,/0/default.jpg}{Sp.~20}. Die Werkausgabe lässt die Tabelle und die Unterschrift aus. Die \href{https://api.digitale-sammlungen.de/iiif/image/v2/bsb10049460_00020/full/1800,/0/default.jpg}{Schlusszeile auf Sp.~32} datiert die Ausgabe auf den 27. Juli 1852. KI-Übertragung und KI-Satz.
\clearpage
